\section{Introduction}
\label{sec:intro}

Harvesting is among the most labour-intensive and time-critical operations in fruit production. The shrinking and increasingly expensive seasonal workforce has made robotic selective harvesting a long-standing goal of agricultural engineering \citep{bac2014harvesting,kootstra2021selective,zhou2022intelligent}. Decades of research have produced capable systems for sweet pepper \citep{arad2020sweetpepper}, kiwifruit \citep{williams2019kiwifruit} and apple \citep{silwal2017apple}. Each, however, is a pipeline of detection, localization, planning and grasping modules engineered for one crop, one decision rule and one scene layout, and a new harvesting objective means a new software project. Vision--language--action (VLA) policies offer a different paradigm. A single network, pretrained on web-scale vision--language data and large robot datasets, maps camera images, proprioception and a natural-language instruction directly to actions \citep{brohan2023rt2,kim2024openvla,black2024pi0,pi05_2025}. For agriculture the appeal is fundamental: \emph{which} fruits to pick, \emph{how many} and \emph{where} to put them become an instruction rather than an engineering effort.

\paragraph{Harvesting is a memory problem} A harvesting robot rarely has in its current view the information it needs for its next decision. Three situations, routine in orchards and greenhouses, make this concrete:
\begin{enumerate}
\item \emph{Foliage occlusion.} Leaves moved by wind or by the manipulator hide ripe fruits that were visible moments earlier \citep{hemming2014detectability}. A fruit that cannot be seen must still be picked.
\item \emph{Quota harvesting.} Grading, sampling and selective thinning require a given number of fruits, so the robot must know how many it has already picked.
\item \emph{Out-of-view containers.} The collection bin is usually outside the working camera's field of view, so the outcome of earlier picks cannot be re-observed.
\end{enumerate}
Forgetting has a direct agronomic cost in each case. A forgotten ripe fruit is lost yield and is left to over-mature on the plant. A forgotten pick leads to over-harvesting that violates the quota. A fruit whose ripeness was misjudged ends as unmarketable produce. The problem is not marginal: in the benchmark introduced here, a leaf-hidden ripe fruit is invisible in a median 78\% of the frames in which it still hangs on the tree, so the decision to pick it must be taken largely from memory.

\paragraph{Time-indexed memory forgets more the longer the robot works} Memory-augmented VLA policies keep a bounded number of past frames, because every frame adds visual tokens and the cost of attention grows quadratically with the number of tokens \citep{vaswani2017attention}. The standard way to fill this budget is to sample the history uniformly in time. We show that this design has a structural flaw for manipulation. The frames that define a pick (the gripper closing on a fruit, the fruit detaching from its stem) span only about a dozen frames, and fewer than 5\% of uniformly sampled memory frames show the gripper within 3~cm of a fruit to be picked. With $B$ frames spread evenly over $t$ elapsed frames, a pick is retained with a probability of roughly $(B-1)|S|/t$, where $|S|$ is the length of its key window. Coverage therefore decays inversely with operating time (Section~\ref{sec:coverage}). The measurements confirm this prediction. At a 16-frame budget, uniform memory misses 43\% of completed picks, can recover the pick count at only 43\% of decisions, and, after 600 frames of operation, loses even the \emph{most recent} pick at 78\% of decisions. These losses are consequential for the policy. On a general manipulation-memory benchmark, removing the memory of the same base VLA reduces success from 66.7\% to 1.7\%, and halving its budget from 32 to 16 frames reduces success from 64.4\% to 26.7\% (Section~\ref{sec:res_general}). What a VLA remembers is a first-order determinant of what it can do.

\paragraph{Why the obvious remedies fall short} \emph{Enlarging the memory} treats the symptom: to keep events, a uniform budget must grow linearly with operating time, and the cost of attention grows quadratically with the budget, which embedded field computers cannot afford. \emph{Better per-frame perception} cannot see what is hidden. A detector without memory loses a fruit as soon as leaves cover it and mistakes a carried fruit for a new target. In our experiments, per-frame detection recovers the correct pick count in only 62\% of episodes. \emph{Learned memory selection or token compression} adds parameters and training signals, may behave differently in training and deployment, and makes it hard to audit what the robot remembered when it failed. What is missing is a \emph{principle} that tells a harvesting robot what is worth remembering.

\paragraph{Key insight: index memory by events, not by time} Our central observation is that in agricultural manipulation, task-relevant information is not spread uniformly over time. It is concentrated at a few physical interaction events. Every pick consists of an approach, a contact and a detachment, and these moments change the state of the task: a fruit leaves the plant and the count increments. They are also physically observable, because the distance between the tool and the nearest fruit still to be picked drops below the fruit radius at contact and jumps at detachment. This mirrors how people remember activity. Humans segment continuous experience into discrete events and preferentially encode event boundaries \citep{zacks2007event}, and episodic memory binds \emph{what} happened to \emph{where} and \emph{when} it happened \citep{tulving1972episodic,clayton1998episodic}. We carry this structure into robot memory.

\paragraph{HarvestMem: event-centric spatial--episodic memory} We introduce \emph{HarvestMem}, a framework that equips a VLA harvesting policy with two complementary memories coupled through one physical signal (Fig.~\ref{fig:overview}).
\begin{itemize}
\item \textbf{Spatial memory} answers \emph{where} the fruits are. A SAM-based fruit distance module (\fdm) segments fruits with the Segment Anything Model \citep{kirillov2023sam} while they are visible and anchors each one in 3D. The anchors persist when foliage hides the fruit and are retired when the fruit is picked.
\item \textbf{Episodic memory} answers \emph{what} happened and \emph{when}. Picking-Event Keyframe Memory (\pekm) detects picks causally from the distance between the tool and the nearest anchored fruit, and reserves slots in the VLA's fixed memory for their contact and detachment keyframes.
\end{itemize}
The design changes \emph{what} the policy remembers, not \emph{how much} or \emph{how}. \pekm\ adds no parameters, no tokens and no inference cost. It contains the base model's uniform memory as the special case $k=0$, and it is applied identically in training and deployment. Because perception and memory communicate only through a scalar distance, the detector can be replaced without retraining the policy, and every memory the policy saw can be reproduced from a logged signal. Phase-aware recovery demonstrations (DART-H) complement the memory by addressing the compounding errors of closed-loop imitation \citep{ross2011dagger}.

\paragraph{Contributions} The main contributions of this work are:
\begin{enumerate}
\item \textbf{Problem formulation and benchmark.} We formulate selective harvesting with a VLA policy as a fixed-budget memory-allocation problem. We release \bench, a procedural ManiSkill3 orchard with a commercial 7-DoF RealMan RM75-B arm and two tasks that cannot be solved from the current view: picking leaf-hidden ripe fruits (FruitOccludedRipe) and meeting a quota with an out-of-view basket (FruitCount). The benchmark includes per-frame ground truth for events, occlusion and perception.
\item \textbf{Event-centric memory with a coverage analysis.} We identify a structural limitation of time-indexed memory: its coverage of interaction events decays inversely with operating time. A simple sampling model predicts the measured coverage to within 0.012. We propose \pekm, which guarantees coverage of the most recent picks independently of episode length. It retains 99.6\% of completed picks against 57.0\% for uniform memory at the same budget, and with 8 frames it outperforms uniform memory with 16, at zero additional cost.
\item \textbf{Spatial--episodic coupling through promptable segmentation.} \fdm\ turns RGB-D images and proprioception into 3D fruit anchors, using visual (depth vacancy) and proprioceptive (grasp) evidence of picking, and reproduces the ground-truth picking events exactly. Its ablations establish that spatial memory acquired before occlusion is a \emph{precondition} for episodic memory: when anchors are built after occlusion, the occluded-task count accuracy falls from 1.00 to 0.00. A robustness analysis turns event memory into an explicit accuracy specification for perception modules.
\item \textbf{Phase-aware recovery demonstrations (DART-H).} Temporally correlated joint noise is injected only in free-motion phases, while clean planner actions are recorded as labels. This targets the covariate shift that limits closed-loop imitation in harvesting.
\item \textbf{A rigorous evaluation protocol.} Evaluation follows a protocol sealed before the formal runs. It uses paired bootstrap with Holm correction and equivalence margins, and audits at every decision that the frames fed to the model match the declared selection rule. Counterfactual controls separate the effect of event information from the mere perturbation of changing frames.
\end{enumerate}

The remainder of the paper is organized as follows. Section~\ref{sec:related} positions the work. Section~\ref{sec:benchmark} describes \bench. Section~\ref{sec:method} presents HarvestMem: the coverage analysis, \pekm, \fdm\ and DART-H. Section~\ref{sec:design} details the experimental design and Section~\ref{sec:results} reports the results. Section~\ref{sec:discussion} discusses implications and limitations for agricultural deployment, and Section~\ref{sec:conclusions} concludes.

\begin{figure}[t]
\centering
\fbox{\parbox{0.92\linewidth}{\centering\vspace{2.2cm}\tbd{Fig.~1 (to be drawn). Left: the two memory needs of harvesting, a ripe fruit hidden by wind-blown leaves and a quota with an out-of-view basket. Middle: time-indexed (uniform) memory misses the short pick events. Right: HarvestMem, with spatial memory (3D fruit anchors from SAM, \emph{where}) and episodic memory (contact/detachment keyframes from PEKM, \emph{what} and \emph{when}) coupled by the tool-to-fruit distance $d_t$ and feeding the VLA.}\vspace{2.2cm}}}
\caption{HarvestMem in a nutshell. Harvesting requires remembering fruits that are no longer visible and picks whose results are out of view. Instead of sampling memory uniformly in time, HarvestMem indexes it by physical events: fruits are anchored in 3D while visible (spatial memory), and the contact and detachment keyframes of picks are reserved in the VLA's fixed memory budget (episodic memory).}
\label{fig:overview}
\end{figure}
